

\section{Basics}

\begin{frame}{Qu'est-ce que Cassandra\,?}

Cassandra, conçu à l'origine par Facebook, maintenant un projet de la fondation Apache,
distribué par la société Datastax.

\begin{redbullet}
  \item Initialement basé sur le système BigTable de Google (stockage orienté colonnes)
   \item A fortement évolué vers un \textit{modèle relationnel étendu}
  \item Un des rares systèmes NoSQL à proposer un typage fort
  \item Un langage de définition de schéma et de requêtes, CQL (pas de jointure)
    \end{redbullet}
   \vfill
   
Construit dès l'origine comme un système \alert{scalable} et \alert{distribué}.

\vfill

\alert{Dans ce cours, nous étudions le modèle de données.}

\end{frame}


\section{le modèle de données}


\begin{frame}
\frametitle{Modèle = relationnel + types complexes}

Cassandra gère des \alert{bases} (\textit{keyspaces}), contenant des \alert{tables} (a.k.a.
\textit{column families}).

\vfill

Une table contient des \alert{lignes} (\textit{rows}) constituées de \alert{valeurs simples}
ou \alert{complexes}.

\vfill

\alert{Perspective A}: c'est du relationnel étendu en rompant la première règle de normalisation (type atomique).

\vfill

\alert{Perspective B}: chaque ligne est un document structuré au sens où nous l'avons vu jusqu'ici.

\vfill

\begin{remblock}{Attention}
Vocabulaire confus et parfois déroutant: \textit{columns}, \textit{column families}, etc.
\end{remblock}

\end{frame}


\begin{frame}
\frametitle{Paires et documents}

La structure de base est la paire (clé, valeur)

\begin{redbullet}
 \item Clé\,: un identifiant
 \item Valeur\,:  atomique (entier, chaîne de caractères), ou complexe (dictionnaire, ensemble, liste)
\end{redbullet}

\vfill

Une ligne (\textit{row}, document) est un identifiant (\textit{row key}) associé à un ensemble de paires (clé, valeur).

\vfill

Les lignes sont \alert{typées} par un schéma, y compris les données imbriquées.

\vfill

Cassandra n'autorise pas (au moins avec CQL) l'insertion de données ne respectant pas le schéma.

\vfill

\alert{Quelques aspects ignorés ici}, le versionnement notamment.
\end{frame}


\begin{frame}
\frametitle{Tables et base de données}

Une \alert{table} (anciennement \textit{column family}) est un ensemble de \textit{rows} conformes
au schéma de la table.

\vfill

Les tables sont \alert{dénormalisées}\,: on peut avoir des attributs dont
la valeur est un type construit (objet ou ensemble).
\vfill

Des notions de colonnes et super-colonnes sont (semble-t-il) en voie de disparition.

\vfill

Une \alert{base} est un ensemble de tables. On l'appelle \textit{keyspace}. Explication? Plus tard.

\end{frame}


\begin{frame}
\frametitle{Conception d'un schéma}

\alert{Cassandra ne propose pas de jointure}. C'est du NoSQL après tout.

\vfill

Il est donc préférable de \alert{regrouper} les données le plus possible dans des lignes
(notion de document autonome, déjà vue).

\vfill


Au pire, on peut représenter les mêmes données sous plusieurs formes.

\begin{redbullet}
  \item Tous les films, avec leurs réalisateurs et acteurs.
  \item Tous les artistes, avec les films qu'ils ont tournés ou dans lesquels ils ont joué.
\end{redbullet}


\vfill

Le regroupement a aussi des inconvénients (redondance, difficulté d'accès aux
données ``profondes''). 
La conception Cassandra s'appuie sur des heuristiques
spécifiques\,: une autre session lui est consacrée.
\end{frame}

\section{Session découverte\,!}


\begin{frame}[fragile]
\frametitle{Application !}

Création d'un \textit{keyspace}.

\begin{minted}{sql}
create keyspace if not exists Movies 
       with replication = { 'class' : 'SimpleStrategy',
               'replication_factor': 3 };
\end{minted}

\vfill

Création d'une table (sans imbrication).


\begin{minted}{sql}
 create table artists (id text, 
                last_name text, first_name text, 
                birth_date int, primary key (id) 
             );
\end{minted}
\end{frame}


\begin{frame}[fragile]
\frametitle{Insertions}

A la SQL

\begin{minted}{sql}
insert into artists (id, last_name, first_name, birth_date)  
      values ('artist1', 'Depardieu', 'Gerard', 1948);
insert into artists (id, last_name, first_name, birth_date)  
          values ('artist2', 'Baye', 'Nathalie', 1948);
insert into artists (id, last_name, first_name)  
          values ('artist3', 'Marceau', 'Sophie');
\end{minted}

\vfill

Avec JSON

\begin{minted}{sql}
    insert into artists JSON '{
         "id": 1,
         "last_name": "Coppola",
         "first_name": "Sofia",
         "birth_date": "1971"
       }' 
\end{minted}
\end{frame}



\begin{frame}[fragile]
\frametitle{Imbrication de structures}

C'est le principe de dénormalisation\,: on regroupe 
les données le plus possible dans des lignes pour éviter les jointures.

\vfill

Une valeur d'attribut  peut correspondre\,:
\begin{redbullet}
	\item à un ensemble de valeurs (non ordonnées)\,: \textbf{SET} 
	\item à une liste de valeurs\,: \textbf{LIST}
	\item à un dictionnaire\,: \textbf{MAP} 
	\item à un nuplet\,: \textbf{TUPLE} 
	\item à une instance d'un type utilisateur
\end{redbullet}

\vfill

Les ensembles doivent rester de taille raisonnable
\end{frame}

\begin{frame}[fragile]
\frametitle{Avec imbrication}

Création d'un type
\begin{minted}{sql}
create type artist (id text,
                    last_name text,
                    first_name text,
                    birth_date int,
                    role text);\end{minted}

\vfill

Table des films, avec imbrication d'un object conforme au type \texttt{artist}.

\begin{minted}{sql}
create table movies (id text,
                     title text,
                     year int,
                     genre text,
                     country text,
                     director frozen<artist>,
                     primary key (id) );
\end{minted}
\end{frame}


\begin{frame}[fragile]
\frametitle{Insertion de documents}

À partir d'un document JSON

\begin{minted}{sql}
insert into movies JSON '{
    "id": "movie:1",
    "title": "Vertigo",
    "year": 1958,
    "genre": "drama",
    "country": "USA",
    "director": {
        "id": "artist:3",
        "last_name": "Hitchcock",
        "first_name": "Alfred",
    "   birth_date": "1899"
    },
}';
\end{minted}

\end{frame}


\begin{frame}[fragile]
\frametitle{Table des films, complète}

On peut avoir des valeurs instances d'un type, en \alert{ensemble} d'instance
d'un type (mot-clé \texttt{set}). Exemple\,: 

\vfill

\begin{minted}{sql}
create table movies (id text,
              title text,
              year int,
              genre text,
              country text,
              director frozen<artist>,
              actors set< frozen<artist>>,
           primary key (id) );
\end{minted}

\end{frame}



\begin{frame}[fragile]
\frametitle{À retenir}

\begin{redbullet}
   \item Cassandra propose un modèle relationnel étendu, basé sur la capacité à imbriquer
     des types complexes dans la définition d'un schéma
   \item  Cassandra a choisi d'imposer un typage fort: toute insertion doit être conforme au schéma;
   \item L'imbrication des constructeurs de type
   rend le modèle comparable aux documents structurés JSON ou XML.
\end{redbullet}

\end{frame}